\chapter{枚乘}

\begin{authorintro}
枚乘（？—前140？），字叔，淮阴（今江苏淮安西南）人。初为吴王刘濞郎中，吴王欲谋反，枚乘上书劝谏，吴王不听，枚乘投奔梁孝王刘武。景帝时，吴王参与六国谋反，乘又上书劝阻，吴王仍然不听，最后兵败身死。枚乘也因此知名于世。七国之乱平定后，景帝召拜乘为弘农都尉。“乘久为大国上宾，与英俊并游，得其所好，不乐郡吏，以病去官”（《汉书·贾邹枚路传》）。乘再度游梁，以辞赋见称。梁孝王死后，乘归淮阴。武帝久闻枚乘之名，即位后，乃以安车蒲轮征乘。因乘时已年老，死于途中。今存赋三篇，其中以《七发》最为著名。另有《谏吴王书》、《重谏吴王书》保存在《汉书》本传中。
\end{authorintro}

\section[七发]{七发}

楚太子\endnote{此文选自《文选》卷三十四。《七发》是汉代大赋中讽喻性较强的作品。赋用主客问答的形式，借设楚太子有疾，而吴客往问之。吴客分别用七件事来启发诱导楚太子，先分别描写音乐、饮食、车马、游宴、田猎、观涛六事，太子皆曰“仆病未能”，最后进“要言妙道”，太子“霍然病已”。作品旨在劝诫贵族集团不要贪欲过度，享乐无时。沉溺于安逸享乐的生活中，是百病之源，这就给予贵族集团的养尊处优、骄奢淫逸不啻提出严重的警告。它用铺张、夸饰的手法来描绘音乐、饮食以及田猎、观涛之事，辞藻华美，结构宏伟。作者善于运用形象的比喻刻画事物，尤其是观涛一段，逼真传神。虽有汉代大赋的铺张扬厉，却无堆垛寡变之嫌。《七发》的出现，标志着汉代散体大赋的形式已基本形成，在赋的发展史上起着转关的作用，且对后代创作产生了巨大的影响，傅毅的《七激》、张衡的《七辩》、曹植的《七启》、陆机的《七征》、张协的《七命》等，在体式上均受其影响，被称作“七体”。}\endnote{楚太子：赋中假设的人物。}有疾，而吴客往问之\endnote{吴客：吴国的宾客，亦为假设的人物。问：问候，探视。}，曰：“伏闻太子玉体不安\endnote{伏闻：犹在下听说。伏，表示谦虚与恭敬。}，亦少间乎\endnote{少间（jiàn建）：稍有好转。}？”太子曰：“惫\endnote{惫：疲惫无力。}，谨谢客\endnote{谨谢客：恭谨地感谢你这位客人。}。”客因称曰\endnote{因：乘机。称：进言。}：“今时天下安宁，四宇和平，太子方富于年\endnote{方富于年：正当年轻的时候。}。意者久耽安乐\endnote{意者：我意料。耽：沉溺，迷恋。}，日夜无极\endnote{无极：没有尽头。}，邪气袭逆\endnote{邪气：淫邪之气。袭逆：侵犯。}，中若结涩\endnote{中：指胸中。结涩（sè色）：即郁结堵塞。按：涩，本为古代车旁用皮革交错结成的障蔽物。此处通“涩”，指气结不畅。}。纷屯澹淡\endnote{纷屯澹淡：昏愦烦闷的样子（用李善注）。}，嘘唏烦酲\endnote{嘘唏：同“歔欷”，叹息呻吟之声。烦酲（chéng成）：心中烦躁，似酒醉未解。}。惕惕怵怵\endnote{惕惕怵怵：戒惧警惕、心神不宁的样子。}，卧不得瞑\endnote{瞑（mián眠）：通“眠”。《文选·嵇康〈养生论〉》“达旦不瞑”，李善注：“瞑，古眠字。”}。虚中重听\endnote{虚中：中气不足。重听：指耳鸣。}，恶闻人声。精神越渫\endnote{越渫（xiè谢）：涣散。}，百病咸生\endnote{咸：都，皆。}。聪明眩矅\endnote{聪明：指耳和眼。眩矅（yào耀）：惑乱昏花。}，悦怒不平\endnote{悦怒不平：喜怒无常。}。久执不废\endnote{执：持而不变，指病缠身。废：止。}，大命乃倾\endnote{大命：指生命。乃倾：将倾覆不保。}。太子岂有是乎？”太子曰：“谨谢客，赖君之力，时时有之，然未至于是也\endnote{“赖君之力”三句：言依赖我的国君的力量，使我享受种种安乐，如客所说，虽时而有这种情况，还不至于像你说得如此严重（用明人张凤翼《文选纂注》说）。}。”客曰：“今夫贵人之子，必宫居而闺处\endnote{宫居而闺处：即住在宫室，处于深闺。}，内有保母\endnote{内：指宫中。保母：照顾太子生活的妇女。}，外有傅父\endnote{外：指朝廷之上。傅父：教育太子的师傅。}，欲交无所\endnote{欲交无所：指太子想出外交游也无机会。}。饮食则温淳甘膬\endnote{温淳：指味道醇厚的食物。甘膬：甜而可口的食物。膬，同“脆”。}，脭醲肥厚\endnote{脭（chéng成）醲肥厚：实即“脭肥醲厚”，指所食精美之肉很肥，所饮之酒很醇厚。吕延济注：“脭，肉之精。”醲，味浓的酒。}。衣裳则杂遝曼暖\endnote{杂遝（tà踏）：众多的样子。曼暖：轻细温暖。}，燂烁热暑\endnote{燂（xún巡）烁：热，炽热。}。虽有金石之坚，犹将销铄而挺解也\endnote{“虽有”二句：言即使身体坚如金石，也会销镕解体。}，况其在筋骨之间乎哉？故曰：纵耳目之欲，恣支体之安者，伤血脉之和\endnote{“纵耳目”三句：言放纵自己的声色享受的欲望，恣意贪图身体的安逸，是有伤血脉调和的。支，同“肢”。}。且夫出舆入辇，命曰蹶痿之机\endnote{“且夫”二句：言出入不步行而乘坐车辇，这正是给腿脚麻痹瘫痪造成机会。命，即名。蹶痿，腿不能走路曰蹶，麻痹曰痿（用吕向注）。机，犹弩机，言触之即发也（用张凤翼说）。按：此二句本于《吕氏春秋·本生》：“出则以车，入则以辇，务以自佚，命曰佁蹷之机。”}。洞房清宫\endnote{洞房：深邃的房屋。清宫：清凉的宫殿。}，命曰寒热之媒\endnote{寒热之媒：寒疾热病的媒介。}。皓齿娥眉\endnote{皓齿娥眉：代指美女。}，命曰伐性之斧\endnote{伐性之斧：戕害生命的斧钺。}。甘脆肥脓\endnote{肥脓：同“肥醲”。注见〔31〕。}，命曰腐肠之药。今太子肤色靡曼\endnote{肤色：指外表。靡曼：纤弱柔美。}，四支委随\endnote{委随：萎弱，软弱。委，通“萎”。}，筋骨挺解，血脉淫濯\endnote{淫濯：谓过度而且大。淫与濯都是大的意思。}，手足堕窳\endnote{堕窳（yǔ雨）：懒散无力。}；越女侍前，齐姬奉后\endnote{“越女”二句：言侍奉身边者皆来自各地的美女。越女、齐姬，指越国与齐国之美女。}，往来游宴\endnote{游宴：亦作“游燕”、“游宴”，游乐宴饮。}，纵恣于曲房隐间之中\endnote{纵恣：任意放肆。曲房：深曲的房屋，犹上文之“洞房”。隐间：隐蔽的秘室。}。此甘餐毒药，戏猛兽之爪牙也\endnote{戏猛兽之爪牙：戏弄猛兽的爪牙。喻危险。}。所从来者至深远\endnote{“所从来者”句：言过这种骄奢淫逸的生活由来已久。}，淹滞永久而不废\endnote{淹滞：淹留沉滞。指沉溺而不能自拔。}，虽令扁鹊治内\endnote{扁鹊：先秦时名医，姓秦名越人，《史记》有《扁鹊仓公列传》。治内：指医治脏腑内的疾病。}，巫咸治外\endnote{巫咸：古神传说中的人名，其时代与职业说法不一。一说黄帝时人，善卜筮。一说唐尧时人，为帝尧的医生。一说殷中宗时的贤人，为神巫。从“巫咸治外”判断，作者似将他视为善治外科疾病的名医。}，尚何及哉！今如太子之病者，独宜世之君子\endnote{独宜：只适合。世之君子：当指后文的“方术之士”，指能“论天下之精微，理万物之是非”者。}，博见强识\endnote{博见强识：见闻广博而记忆力强。识，音义同“志”。}，承间语事\endnote{承间：乘空闲之时。犹伺机。语事：谈所闻所见之事。}，变度易意\endnote{变度：改变平度（平时）的生活方式。易意：改变（太子的）意念。}，常无离侧，以为羽翼\endnote{羽翼：指辅佐。}。淹沉之乐\endnote{淹沉之乐：指沉溺心志的娱乐。}，浩唐之心\endnote{浩唐之心：浩荡、放荡之心。按：浩荡亦含荒唐义。《楚辞·离骚》：“怨灵修之浩荡兮，终不察夫民心。”姜亮夫校注：“浩荡，犹今言荒唐耳，一声之转也。”}，遁佚之志\endnote{遁佚：放纵，淫佚。}，其奚由至哉\endnote{其：指代上文的“淹沉之乐”等。奚：何。}！”太子曰：“诺\endnote{诺：好。应答之辞。}。病已，请事此言\endnote{“病已”二句：言等我的病好了，一定照你说的去做。}。”

客曰：“今太子之病，可无药石针刺灸疗而已\endnote{“可无”句：言不用药物针灸便可治好病。}，可以要言妙道说而去也\endnote{要言：中肯的至理名言。妙道：精妙的道理。说（shuì税）：说服，劝诱。去：指除去疾病。}，不欲闻之乎？”太子曰：“仆愿闻之。”

客曰：“龙门之桐\endnote{龙门之桐：产于龙门山的桐木。龙门山在山西河津和陕西韩城之间。桐木最宜制琴。}，高百尺而无枝，中郁结之轮菌\endnote{中：指树干。郁结：积聚。轮菌：纹理盘曲的样子。}，根扶疏以分离\endnote{扶疏：指根向四面分布伸展。分离：指大根小根交相分开。}。上有千仞之峰，下临百丈之溪，湍流溯波，又澹淡之\endnote{“湍流”二句：言树根经受过急流逆波的冲击、摇荡。澹（yǎn演）淡：水波动荡的样子，此处用作水波对树根的激荡。}。其根半死半生，冬则烈风、漂霰、飞雪之所激也\endnote{漂霰：同“飘霰”，飘落的雪珠。激：激荡。}，夏则雷霆、霹雳之所感也\endnote{感：感动，感染。按：龙门之桐，其所以能成为制琴之佳木，在于受自然界各种音响的感染。故“感”字训为感动、感染可通。有的注本认为“感”是“撼”的假借字，解作“震撼”，亦可通。}。朝则鹂黄、鶟鶗鸣焉\endnote{鹂黄：鸟名，即黄鹂，又名黄莺，亦称鸧鹒。鶟鶗（hàn dàn汉旦）：亦作“鹖旦”，鸟名，似鸡，冬无毛，昼夜常鸣。}，暮则羁雌、迷鸟宿焉\endnote{羁雌：失偶的雌鸟。迷鸟：迷失方向的飞鸟。}。独鹄晨号乎其上\endnote{独鹄：孤独的黄鹄。黄鹄，俗名天鹅。}，鹍鸡哀鸣翔乎其下\endnote{鹍（kūn昆）鸡：鸟名，似鹤，黄白色。}。于是背秋涉冬\endnote{背秋涉冬：秋天过去进入冬天。意谓不知经过了多少个寒暑。此句写桐树的生长过程。}，使琴挚斫斩以为琴\endnote{琴挚：人名，指师挚，古代有名的琴师。李善注引郑玄说：“师挚，鲁太师（乐官名）也，以其工琴，谓之琴挚。”斫斩以为琴：把桐木砍下来截断制成琴。}，野茧之丝以为弦\endnote{“野茧”句：用野蚕所吐之丝做琴弦。}，孤子之钩以为隐\endnote{孤子之钩：指孤儿衣带上的钩。隐：琴上饰物。李善注引桓子《新论》说：“琴隐长四十五分。”按：孤儿命苦，琴声以悲为美，故用之以增其悲苦。}，九寡之珥以为约\endnote{九寡：生有九个儿子的寡妇。李善注引《列女传》说：“鲁之母师（女琴师），九子之寡母也。不幸早少夫，独与九子居。”珥：珠玉做的耳饰，也叫瑱、珰。约：琴徽，为琴上指示音阶的标志。}。使师堂操《畅》\endnote{师堂：一称师襄，字子京，古代的乐师。孔子曾学鼓琴于师襄子（见《韩诗外传》卷五第七章）。操：弹奏。《畅》：尧时的琴曲。又称《神人畅》（用胡绍煐引陈旸《乐书》说）。}，伯子牙为之歌\endnote{伯子牙：即伯牙，古之善琴者。事见《吕氏春秋·本味篇》：“伯牙鼓琴，钟子斯听之，方鼓琴而志在太山，钟子期曰：‘善哉乎鼓琴，巍巍乎若太山。’少选之间，而志在流水，钟子期又曰：‘善哉乎鼓琴，汤汤乎若流水。’钟子期死，伯牙终身不复鼓琴。”}。歌曰：‘麦秀蔪兮雉朝飞，向虚壑兮背槁槐\endnote{“麦秀”二句：言当麦子抽穗生芒时，野鸡在早晨飞去，野鸡离开枯槐又向空谷飞去。麦秀，指麦穗。蔪（jiān尖），指麦芒。虚壑，空虚的沟壑。背，背向，引申为离去。槁槐，即枯槐。}，依绝区兮临回溪\endnote{依：靠近。绝区：险绝之地。回溪：曲折的溪流。此句是上句的补充，接写雉鸟飞向依山临水之处。}。’飞鸟闻之，翕翼而不能去\endnote{翕（xī吸）翼：收敛翅膀。去：离开，飞去。}；野兽闻之，垂耳而不能行\endnote{垂耳：两耳下垂。形容驯服的样子。}；蚑蟜蝼蚁闻之\endnote{蚑蟜（qí jiǎo其矫）：指爬行的小虫。蚑，为蟢蛛。蟜为一种毒虫。}，拄喙而不能前\endnote{拄喙：把嘴支在地上。}。此亦天下之至悲也\endnote{至悲：此指最美的音乐。按：古代以悲为美。}。太子能强起听之乎？”太子曰：“仆病，未能也。”

客曰：“犗牛之腴\endnote{犗（chú除）牛：小牛。腴：腹部肥肉。}，菜以笋蒲\endnote{菜以笋蒲：用竹笋和蒲菜做犗牛的配菜。菜，作动词用，“菜以”犹“配以”。蒲，即香蒲，又名甘蒲。丛生水际，根茎可食，叶可制席、扇等。笋与蒲是我国古代的重要蔬菜。《诗·大雅·韩奕》：“其蔌（蔬菜）维何？维笋及蒲。”是其证。}；肥狗之和\endnote{肥狗之和：用肥美的狗肉调成的肉羹。和，指调和成羹。}，冒以山肤\endnote{冒：覆盖。山肤：即石耳，是地衣类植物，扁平如叶，圆形，呈灰黑色，附着于岩石，可以采食，被称作“山珍”。此句言在狗肉羹的表面撒上石耳调成肉菜混合的羹。}。楚苗之食\endnote{楚苗之食：用楚苗山出产的大米所做之饭（用李善注）。}，安胡之饭\endnote{安胡之饭：用雕胡米做成的饭。安胡，即雕胡，又称菰米。}，抟之不解，一啜而散\endnote{“抟之”二句：言优质米饭很粘，用手团紧米团不会散开，一吃到口里便散开了。啜（chuò辍），食，饮。}。于是使伊尹煎熬\endnote{伊尹：商汤大臣，名伊，一名挚，尹是官名。传说他善于烹饪。}，易牙调和\endnote{易牙：春秋时人，又称狄牙、雍巫，为齐桓公的宠臣，长于调味，善逢迎，传说曾烹其子为羹以献桓公。事见《左传·僖公十七年》、《战国策·魏策》、《史记·齐太公世家》。调和：指调和五味。}，熊蹯之胹\endnote{熊蹯（fán凡）：熊掌。胹（ér而）：烂熟的肉。此句言将熊掌烹得烂熟。}，勺药之酱\endnote{勺药之酱：指五味调和的汤汁。据清代沈钦韩《汉书疏证》说，勺药，应读为“酌略”，指调和五味。酱，汤汁。}，薄耆之炙\endnote{薄耆之炙：把兽类的里脊肉切成薄片加以烧烤。耆，指兽脊之肉。}，鲜鲤之鲙\endnote{鲙（kuài快）：鱼肉片。}，秋黄之苏\endnote{秋黄之苏：秋天叶子变黄的紫苏草。苏，即紫苏，又名桂荏。一年生草本植物，茎方形，花淡紫色，种子可榨油，嫩叶可以吃。叶、茎和种子均可入药。}，白露之茹\endnote{白露之茹：白露以后的蔬菜。茹，指蔬菜。有些蔬菜经霜露后变甜，适于食用。}。兰英之酒\endnote{兰英之酒：用兰花浸泡的酒。取其芳香。}，酌以涤口\endnote{酌：饮。涤口：漱口。}，山梁之餐\endnote{山梁：指野鸡。《论语·乡党》：“山梁雌雉，时哉时哉。”孔子的本义感叹山桥上的雌雉得其时而人不得时。后代遂以山梁指代雉，已非山梁之本意。}，豢豹之胎\endnote{豢（huàn宦）豹：被人畜养的豹。胎：指未出生的小豹。豹胎，在古代是一种美食。}。小饭大歠，如汤沃雪\endnote{“小饭”二句：言不管是小吃还是大饮，都像沸水浇在雪上一般，入口便化，非常爽口。歠（chuò啜），饮，指喝汤。沃，浇灌。}。此亦天下之至美也。太子能强起尝之乎？”太子曰：“仆病，未能也。”

客曰：“钟岱之牡\endnote{钟、岱：二地名，古时属赵国，两地以产马著名。钟，即钟山，又称阴山。岱，当作“代”，即今之山西代县。牡，指雄马。}，齿至之车\endnote{齿至：指马的牙齿长满，标志马的筋力成熟。至，达到极点之义。齿至之车指用成年之马所驾之车。}，前似飞鸟，后类距虚\endnote{“前似”二句：言驾车的骏马在前者像飞凫骏马，在后者像形似骡马的距虚兽。飞鸟，当作“飞凫”，骏马名（用清人林茂春《文选补注》引《齐民要术》说）。距虚：传说中的兽名。《吕氏春秋·不广》：“北方有兽，名曰蹶，鼠前而兔后，趋则跆（跌倒），走则颠，常为蛩蛩距虚取甘草以与之。蹶有患害也，蛩蛩距虚必负而走。”距虚，一作“驫駏”（见《尔雅》），形似骡而小。}。穱麦服处\endnote{穱麦：早熟的麦子。服处：饲养。}，躁中烦外\endnote{躁中烦外：指马的内火盛而外表烦躁。按：以穱麦喂马则马的内火盛，内火上升则烦躁，烦躁则亟思奔驰。}。羁坚辔\endnote{羁坚辔：系上坚固的马缰绳。}，附易路\endnote{附易路：依附平坦的道路。}。于是伯乐相其前后\endnote{伯乐：秦穆公时人，善于相马。}，王良、造父为之御\endnote{王良：春秋时代最善于驾车的人，曾为赵简子驾车。造父：周穆王的驾车者。据《穆天子传》说，他曾驾八骏载穆天子西游。御：驾车。}，秦缺、楼季为之右\endnote{秦缺、楼季：皆古之善走者（用李善说）。一说为古代的勇士。右：指车右。}。此两人者\endnote{两人：李善注谓指秦缺、楼季二人。}，马佚能止之\endnote{马佚：指马受惊逃逸。}，车覆能起之\endnote{车覆：翻车。起：扶起。}。于是使射千镒之重\endnote{射：射覆，古时的一种猜物游戏，此指打赌。千镒（yì意）之重：以千镒重金作赌注。镒，古代重量单位，一镒为二十四两。}，争千里之逐\endnote{千里之逐：在骏马的千里追逐中决胜负。}。此亦天下之至骏也\endnote{至骏：最好的马。}。太子能强起乘之乎？”太子曰：“仆病，未能也。”

客曰：“既登景夷之台\endnote{景夷之台：景夷，台名，据王文濡《古文辞类纂音注》说，景夷台即京台，一作“荆台”、“强台”，在今湖北监利北。又据东汉边让《章华台赋序》：“楚灵王既游云梦之泽，息于荆台之上，前方淮（水名）之水，左洞庭之波，右顾彭蠡之隩，南眺巫山之阿，延目广望，骋观终日，顾谓左史倚相曰：‘盛哉此乐，可以遗老而忘死也。’于是遂作章华之台。”看来筑章华台比荆台晚。又《史记·楚世家》云：“（楚灵王）七年，就章华台，下令内亡人实之。”史记《集解》引杜预注：“南郡华容县（今属湖南）有台，在城内。”}，南望荆山\endnote{荆山：山名，在今湖北南漳西。按：景夷台在湖北监利，荆山在其大西北六七百里之处，不得言南望。王文濡认为“荆山”当为“猎山”之误，猎山当在湖南华容县境，似可从。又刘向《说苑·正谏》：“楚昭王欲之荆台游，司马子綦进谏曰：‘荆台之游，左洞庭之波，右彭蠡之水，南望猎山，下临方淮。’”可反证“南望荆山”，当为“南望猎山”之误。}，北望汝海\endnote{汝海：即汝水，源出河南鲁山县大盂山，东流而注入淮河。李善注：“汝称海，大言之也。”}，左江右湖\endnote{江：指长江。湖：指洞庭湖。}，其乐无有\endnote{无有：无有过此的省文。}。于是使博辩之士\endnote{博辩之士：博学善辩之人。}，原本山川\endnote{原本山川：陈说山川的本原。}，极命草木\endnote{极命草木：把草木的名字全部说出。极，尽。命，名。}，比物属事\endnote{比物属事：把各种事物排比连缀起来。比，排列。物与事，指山川、草木。属，连缀。}，离辞连类\endnote{离辞连类：把事物按类附丽成文辞。离，通“丽”，附丽，附着。}；浮游览观，乃下置酒于虞怀之宫\endnote{虞怀之宫：指可以赏心娱怀的宫室。虞，通“娱”。}。连廊四注\endnote{连廊四注：宫室的回廊四面相连。注，连、通之意。}；台城层构\endnote{台城：指台上有建筑物。层构：一层层地建造。}，纷纭玄绿\endnote{纷纭玄绿：指建筑物上涂有黑色绿色缤纷夺目。}，辇道邪交\endnote{辇道：专供帝王使用的车道。邪交：纵横交错。}，黄池纡曲\endnote{黄池：即“潢池”，又称“湟池”、“隍池”。围绕城墙的积水池。今称护城河。纡曲：曲折。}。溷章白鹭\endnote{溷（hún混）章：李善注：“溷章，鸟名，未详。”吴小如《枚乘〈七发〉李善注订补》：“溷章，毛色不一之鸟也，故与‘白鹭’为对文。”}，孔鸟鹍鹄\endnote{孔鸟：即孔雀。鹍（kūn昆）鹄：一种大的水鸟。薛综《西京赋》（张衡著）注：“鹍，大鸟也。”鹍鹄之“鹍”，或取其大之意。}，鹓雏鷩鸂\endnote{鹓雏（yuān chú鸳除）：传说中与鸾凤同类的鸟。鷩鸂（jiāo jīng交精）：即池鹭。李时珍《本草纲目·禽一·鷩鸂》：“鷩鸂大如凫、鹜而高脚，似鸡，长喙好啄，其顶有红毛如冠，翠鬣碧斑，丹嘴青胫，养之可玩。”}，翠鬣紫缨\endnote{翠鬣（liè列）：翠绿色的冠毛。紫缨：紫色的颈毛。此句是形容鹓雏、鷩鸂等鸟的毛羽的。}。螭龙德牧\endnote{螭（chī吃）龙：本指雌龙与雄龙。但此处描写鸟，与龙非一类，乃以龙之雌雄借言鸟。德：指凤凰头上的花纹。《山海经·中山经》：“五采而文，名曰凤皇，首文曰德。”牧：本指牛腹的花纹像牧字。见《尔雅·释兽》：“黑腹，牧。”此代指鸟腹上的花纹。}，邕邕群鸣\endnote{邕邕：群鸟和鸣的声音。}。阳鱼腾跃\endnote{阳鱼：即鱼类。李善注引曾子曰：“鸟鱼皆生于阴，而属于阳。”故称鱼为阳鱼。}，奋翼振鳞\endnote{奋翼：指鱼摇动鱼鳍。鳍如鱼身上的翅膀，故称翼。振鳞：振动鳞甲。此句写鱼游水中之状。}。漃漻菭蓼\endnote{漃漻菭（chóu仇）蓼：言水清净之处，则生菭蓼二草。菭，即蕏草。蓼，水草。（用李善注）}，蔓草芳苓\endnote{蔓草：长有长茎的草。芳苓：香草名。《汉书·扬雄传上》：“愍吾累之众芬兮，飓烨烨之芳苓。”颜师古注：“苓，香草名，音零。”}。女桑河柳\endnote{女桑：柔嫩的小桑树。河柳：即柽柳。也称观音柳、西河柳、三春柳、红柳等。为落叶亚乔木，赤皮，枝细长，多下垂。}，素叶紫茎\endnote{素叶：指女桑。素，色单纯。紫茎：指河柳。因其皮赤，近紫，故言。}。苗松豫章\endnote{苗松：李善注说：“苗松，未详。一曰苗山之松。”按：苗山即苗岭，乃南岭之正干，自云南入贵州境，横贯南部，以入湖南。豫章即樟树。}，条上造天\endnote{条上造天：枝条上达天空。指苗松与樟树而言。}。梧桐并闾，极望成林\endnote{“梧桐”二句：言梧桐和棕榈众多，极目望去，已形成树林。并闾，一作栟榈，即棕榈树，为常绿乔木。}。众芳芬郁\endnote{众芳：指众多的草木。芬郁：香气浓郁。}，乱于五风\endnote{五风：五方之风。《文选》李周翰注：“五风，宫商角徵羽之风也。”古以宫、商、角、徵、羽配东、西、南、北、中五方。}。从容猗靡\endnote{从容：指树木在风中仍呈现从容的姿态。猗靡：指树木被风吹得东歪西倒。猗靡，犹“披靡”。}，消息阳阴\endnote{消息阳阴：指树叶被风吹动，阳阴（正反）两面时隐时现。消，灭。息，生。消息，即忽生忽灭，忽有忽无，引申为或隐或现。}。列坐纵酒\endnote{列坐：按次序排开座位。纵酒：纵情饮酒。}，荡乐娱心\endnote{荡乐娱心：放荡于乐舞之中以使人们内心高兴。}。景春佐酒\endnote{景春：战国时的纵横家。《孟子·滕文公下》：“景春曰：‘公孙衍张仪，岂不诚大丈夫哉，一怒而诸侯惧，安居而天下熄。’孟子曰：‘是焉得为大丈夫乎！’”赵岐注：“景春，孟子时人，为纵横之术者。”佐酒：劝饮。}，杜连理音\endnote{杜连：春秋时善琴者。传说为伯牙之师。《文选》刘良注：“杜连即田连，善鼓琴者。”理音：调理音调。实即奏乐。}。滋味杂陈\endnote{滋味：代指各种美味佳肴。杂陈：错杂陈列。}，肴糅错该\endnote{肴糅：各种肉肴。《一切经音义》四引《通俗文》：“肴杂曰糅。”错该：错杂地陈列于前。该，具备。}。练色娱目\endnote{练色：精选美色。练，精简、选择之意。}，流声悦耳\endnote{流声：指当时人们所爱好的淫乐，如郑卫之声等。}。于是乃发《激楚》之结风\endnote{发：与后文之“扬”相对成文，指引声扬喉歌唱。激楚：歌曲名。为激切昂扬的楚歌。结风：指歌曲结尾的馀声、拖音。}，扬郑卫之皓乐\endnote{郑卫：周代分封的两个诸侯国，春秋时郑卫之声代表当时的新乐，与古乐相对抗，被喜爱古乐的人目为淫邪之声，故孔子说：“郑声淫，放郑声。”皓乐：美妙的音乐。}。使先施、征舒、阳文、段干、吴娃、闾娵、傅予之徒\endnote{先施：即美女西施。征舒：李善注谓指征舒之母夏姬。按：夏姬乃郑穆公之女，嫁夏御叔（陈定公明孙子）为内子（妻子），御叔与夏姬生一子，名征舒，字子南。李善既认为先施至傅予七人皆为美女，而征舒为男性，故认为征舒指征舒之母夏姬，夏姬确实是美丽无比，且妖冶淫荡。张云璈《选学胶言》提出不得以夏姬为征舒，亦有一定道理。张云璈又说：“窃谓此七人中，必是杂举男女之美者，以侈陈游宴之乐。”此说近是。征舒在此可作为美男子理解。阳文：相传为古代楚国的美女。《淮南子·修务训》：“曼颊皓齿，形夸骨佳，不待脂粉芳泽而性可说（悦）者，西施、阳文也。”高诱注：“西施、阳文，古之好女。”段干：古代美女名。吴娃：吴国的美女。吴俗以美女为娃。闾娵（zōu邹）：李善注引韦昭《汉书注》：“梁王魏婴之美人。”傅予：古美女名。}，杂裾垂髾\endnote{杂裾：穿着五彩相配的衣裙。杂，同“襍”，《说文》：“襍，五采相合。”垂髾（shāo梢）：垂着燕尾形的发髻。}，目窕心与\endnote{目窕心与：用目光挑逗传情而心中暗暗相许。窕，同“挑”。}；揄流波，杂杜若\endnote{“揄流波”二句：李善注：“言引流波以自洁，杂杜若以为芳。”五臣注《文选》张铣注：“谓目之光若引水之波；杜若香草也，其衣之香与此相杂。”揄，引。流波，指流水与目光皆可通，两说可并存。}，蒙清尘\endnote{蒙清尘：形容美女衣薄如雾縠，看去像蒙上一层薄薄的细尘。}，被兰泽\endnote{被兰泽：言美女发肤上覆盖着兰膏的香泽。被，同“披”，此处义近“蒙”字。}，宴服而御\endnote{宴服：同“嬿服”，即闲居之服，今称便服。御：进御，入侍。李善注引《尚书大传》：“古者后夫人，至于房中，释朝服，袭嬿服，入御于君也。”}。此亦天下之靡丽皓侈广博之乐也\endnote{靡丽：指声而言则为淫靡美妙；指色而言则为姿容艳丽。此句概括以上声色之乐，则应兼指二者。皓侈：鲜明而盛大。指声色之乐的排场。广博：无穷无尽。}。太子能强起游乎？”太子曰：“仆病，未能也。”

客曰：“将为太子驯骐骥之马\endnote{驯：驯服。骐骥：骏马。}，驾飞軨之舆\endnote{飞軨（líng零）：轻便的猎车。车上有窗。《尚书大传》卷二：“未命为士，车不得有飞軨。”郑玄注：“如今窗车也。”}，乘牡骏之乘\endnote{“乘牡骏”句：乘坐壮健的骏马驾的车。前一个“乘”字是动词，后一“乘”字是名词。四马驾一车称一乘。牡，当作“壮”，形近而误（用朱珔《文选集释》说）。}，右夏服之劲箭\endnote{“右夏服”句：右手拿着夏后氏箭袋中强劲的箭。夏，指夏后氏。服：“箙”的假借字。李善注引服虔说：“‘服’，盛箭器也。夏后氏之良弓，名‘繁弱’。其矢亦良。”}，左乌号之雕弓\endnote{乌号：良弓名，相传为黄帝所使用。雕弓：绘饰有花纹图案的弓。}。游涉乎云林\endnote{云林：指云梦之林（用李善注）。云梦为古代楚国的大泽。}，周驰乎兰泽\endnote{“周驰”句：围绕着出产兰草的大泽奔驰。}，弭节乎江浔\endnote{弭节：驻节；停车。节，为车行的节度。江浔：江边。}。掩青薠\endnote{掩青薠：休息于薠草之上。掩，息（用李善注）。}，游清风\endnote{游：五臣注本《文选》作“溯”，义较优。溯清风即迎着清风。}，陶阳气\endnote{陶：畅；舒展。阳气：指春天的气息。}，荡春心\endnote{荡春心：荡漾着春天愉快的心情。}。逐狡兽\endnote{逐狡兽：追逐狡猾不易捕捉的野兽。}，集轻禽\endnote{集轻禽：集中弓箭攒射轻捷善飞的禽鸟。}。于是极犬马之才\endnote{“极犬马”句：极尽猎犬和奔马的技能。}，困野兽之足\endnote{“困野兽”句：使野兽足力困乏，以便捕捉。}，穷相御之智巧\endnote{穷：穷尽。相御：相，指射猎时的向导。《论语集解》引马融说：“相，导也。”御，指驾车者。}。恐虎豹，慴鸷鸟\endnote{“恐虎豹”二句：使虎豹感到恐惧，使凶悍的鸟类慑服。}。逐马鸣镳\endnote{逐马：奔马，指追逐野兽之马。鸣镳（biāo标）：马口所衔之铁发出响声。一说鸣镳谓銮铃（车上的铃）与马勒旁的横铁相撞击发出的声音。}，鱼跨麋角\endnote{鱼跨：李善注谓“跨度鱼也”。按：跨度鱼三字颇为费解。清代学者多未解其义。愚以为“鱼跨”乃鱼跃而跨，是勇士捕捉野兽的动作。麋角：即“角麋”之倒文，即捕取麋鹿。此句言鱼跃而跨在麋鹿之上而捕捉麋鹿。}；履游麕兔\endnote{履游：与下文之“蹈践”同义，同指用脚践踏。麕（jūn君）：即麞子。}，蹈践麖鹿\endnote{麖（jīng京）：大鹿。}。汗流沫坠\endnote{汗流沫坠：指奔驰的猎马汗流满身口沫下坠。}，冤伏陵窘\endnote{冤伏：指禽兽四下逃匿。按：冤，本义为屈缩，引申为龟缩躲藏。伏亦藏也。陵窘：窘迫。李善注：“陵，犹促也。《说文》曰：‘窘，迫也。’”}，无创而死者，固足充后乘矣\endnote{“无创”二句：言禽兽未受伤而死的，足以装满侍从的车辆。后乘，指在天子或诸侯的车后侍从的车辆。}。此校猎之至壮也\endnote{校猎之至壮：打猎的壮观景象。}。太子能强起游乎？”太子曰：“仆病，未能也。”然阳气见于眉宇之间\endnote{阳气：喜色。见：同“现”，呈现。眉宇之间：指眉额之间。}，侵淫而上\endnote{侵淫：渐进的样子。}，几满大宅\endnote{几满大宅：几乎布满整个面部。大宅，指面部，因为五官的居宅，故称大宅。}。

客见太子有悦色，遂推而进之曰：“冥火薄天\endnote{冥火：指夜间的猎火，用以驱赶野兽。薄天：迫天，近天。}，兵车雷运\endnote{雷运：车辆运行，其声如雷。}，旍旗偃蹇\endnote{旍（jīng京）旗：旌旗。偃蹇：高耸的样子。}，羽旄肃纷\endnote{羽旄：用鸟羽和牛尾装饰的旌旗。肃纷：整齐而盛多的样子。}。驰骋角逐\endnote{角逐：竞相追逐。}，慕味争先\endnote{慕味争先：为追求美味而争先恐后。}。徼墨广博\endnote{徼墨广博：因打猎而焚烧的田野边界广阔。徼，边界。墨，烧田。古人为捕捉野兽而纵火除地，叫烧田，烧后土黑，故称“墨”。}，观望之有圻\endnote{圻（yín银）：同“垠”。边际。此句言烧田广博，远远望去只能看到边际。}。纯粹全牺\endnote{纯粹：指禽兽的毛色纯一。全牺：野兽的身体完整叫全，色纯叫牺。}，献之公门\endnote{公门：指诸侯之门。按：古代祭祀要用毛色纯粹的牺牲。献之公门，用于祭祀。}。”太子曰：“善，愿复闻之。”

客曰：“未既\endnote{未既：还没有完。指客谈狩猎之乐的话还没说完。}，于是榛林深泽\endnote{榛林：丛林。深泽：深广的沼泽。}，烟云暗莫\endnote{暗莫：同“暗漠”。昏暗不明的样子。}，兕虎并作\endnote{兕（sì寺）虎：独角野牛与老虎。并作：并起，一块出来。}。毅武孔猛\endnote{毅武：刚强勇武。指打猎的勇士而言。孔猛：非常猛烈。孔，极，甚。}，袒裼身薄\endnote{袒裼（xī夕）：袒露着身子。身薄：以身靠近野兽。指与野兽相搏斗。}，白刃硙硙\endnote{硙（ái挨）硙：同“皑皑”。雪白的样子。此处形容白刃的霜锋。}，矛戟交错。收获掌功\endnote{收获掌功：按猎取野兽之多少决定功劳。掌，主持，掌管。引申为决定。}，赏赐金帛；掩薠肆若\endnote{掩薠：覆盖（铺上）薠草。肆若：陈列杜若（香草）。此句言在薠草杜若上铺设席位。}，为牧人席\endnote{牧人：官名，专管在田野中牧养牲畜，此指参与打猎的官员与侍从。席：坐席。从后文看，似为饮宴之席。}。旨酒佳肴，羞炰脍炙\endnote{羞：同“馐”，有滋味的食物。炰（páo咆）：同急火烹的食物。脍：细切的肉片。炙：烧烤。}，以御宾客\endnote{御：供给，享用。以御宾客，即以美酒佳肴供给宾客享用。}。涌触并起，动心惊耳\endnote{“涌触”二句：把酒注满杯，一起站立欢呼畅饮，欢声惊心动耳。按：涌触，五臣注本《文选》作“涌觞”，繁体的“觸”与“觴”形近，当作“觞”。觞为酒杯。}。诚必不悔\endnote{诚必：诚，指忠诚，诚信；必，通“毕”指已答应别人的事一定做完。即善始善终。诚必二字当连读。《吕氏春秋·论威》篇：“又况乎万乘之国而有所诚必乎？”不悔：不反悔。}，决绝以诺\endnote{决绝：十分坚决；十分肯定。以诺：同“已诺”，指已经应诺之事。此句言已答应之事态度坚定。}；贞信之色，形于金石\endnote{“贞信”二句：言众人的忠贞诚信之心，都借音乐表现出来。李善注引《孔子家语》说：“孔子曰：‘夫钟鼓之音，忧而击之则悲，喜而击之则乐，故志诚感之，通于金石，而况人乎哉？’”}；高歌陈唱，万岁无斁\endnote{“高歌”二句：言众人纵情歌唱，永远不会厌倦。陈唱，列队而唱。斁（yì义），厌倦。}。此真太子之所喜也，能强起而游乎？”太子曰：“仆甚愿从，直恐为诸大夫累耳\endnote{直恐：只恐怕。累：累赘。}。”然而有起色矣\endnote{起色：跃跃欲试的样子。}。

客曰：“将以八月之望\endnote{望：夏历的每月十五曰望。潮水涨潮之高峰，在夏历的十五日左右，此与月球引力有关。}，与诸侯远方交游兄弟，并往观涛乎广陵之曲江\endnote{广陵：即今江苏扬州。曲江：当在扬州城外瓜洲以北。按：广陵之曲江，亦说指浙江之钱塘江，历史上聚讼纷纭。自北魏郦道元将曲江观涛系在浙江之后，开始产生分歧，直至清代犹未解决。清初的毛奇龄、朱彝尊、阎若璩，均将广陵之涛，移于钱塘。汪中、俞思谦等人，则力主扬州说。连清代的乾隆皇帝也参与了讨论。他在《过瓜洲镇》诗“广陵潮昔有明征”句自加注说：“《蔡宽夫诗话》：‘润江大江本与扬子桥对，瓜洲乃江中一洲耳。’李绅诗：‘扬州郭里见潮生。’自瓜洲以闸为限，潮遂不至扬州。夫唐时扬州尚见潮，何况于汉。著作家迁就枚乘赋语，以曲江观涛为是杭非扬，且谓广陵旧治甚大，钱塘当在所辖之内，此与刻舟胶柱之见何异！既为之辨，复注大意于此。”（转引自人民文学出版社《扬州历代诗词》第三册115页）}。至则未见涛之形也。徒观水力之所到，则恤然足以骇矣\endnote{恤（xù续）然：惊恐的样子。骇：惊怕。}。观其所驾轶者\endnote{驾轶：凌驾超越。指水势而言。}，所擢拔者\endnote{擢拔：水势高耸突起的样子。}，所扬汩者\endnote{扬汩（gǔ古）：扬波，激荡。}，所温汾者\endnote{温汾：水回旋、结聚的样子。}，所涤汔者\endnote{涤汔（qì迄）：冲刷涤荡。}，虽有心略辞给\endnote{心略：心计，心智。李周翰注：“心略，心计也。”辞给（jǐ己）：言辞敏捷。}，固未能缕形其所由然也\endnote{“固未能”句：言还不能细致地描绘出江涛从始至终的种种姿态。缕形，雕镂其形状，即细致描绘。由然，原委，来由。}。怳兮惚兮\endnote{怳兮惚兮：义同“怳惚”，一作“恍惚”，此指江涛浩渺无际，看不真切。}，聊兮栗兮\endnote{聊兮栗兮：使人惊惧战栗。李善注：“聊栗，恐惧之貌。”}，混汩汩兮\endnote{混汩汩兮：许多波涛会合所发出的声音。混，读滚，指波涛翻滚不尽。汩汩，波涛声。}；忽兮慌兮\endnote{忽兮慌兮：义同上文“怳兮惚兮”。}，俶兮傥兮\endnote{俶（tì替）兮傥（tǎng淌）兮：义同“倜傥”，卓异，突出。此处形容浪涛突起非同一般。}，浩瀇瀁兮\endnote{浩瀇瀁（wǎng yǎng往养）兮：形容水势浩大深广。瀇瀁，音义皆同“汪洋”。}，慌旷旷兮\endnote{慌旷旷兮：形容水势空阔望不到边。慌，义同“恍”。}。秉意乎南山\endnote{秉意：执意，集中注意力。南山：泛指曲江以上之山，今南京，镇江一带的山脉。}，通望乎东海\endnote{“通望”句：言一直望到东海。即目送江涛入海。}；虹洞兮苍天\endnote{“虹洞”句：言江涛汹涌，与天相连。虹洞，同“澒洞”，水势汹涌。}，极虑乎崖涘\endnote{“极虑”句：竭尽心思欲看到江涛的边缘。崖涘，同“涯涘”，边际。}。流揽无穷\endnote{流揽：一作“流览”，周游观览。}，归神日母\endnote{“归神”句：把注意力集中到东方日出之处。日母，指太阳。}。汩乘流而下降兮\endnote{汩：水流迅急的样子。乘流而下降：指潮头顺着江流向下游流去。}，或不知其所止\endnote{“或不”句：言有的潮头不知止于何处。}。或纷纭其流折兮\endnote{“或纷纭”句：李善注：“言众浪纷纭，其流曲折。”}，忽缪往而不来\endnote{“忽缪往”句：言潮头忽然纠缠在一起一去而不返。缪，纠结。}。临朱汜而远逝兮\endnote{朱汜：李善注谓地名。刘良注谓南方水涯。未知谁是，姑且两说并存。远逝：指浪潮远去。}，如虚烦而益怠\endnote{“中虚烦”句：言因涛逝去内心感到空虚烦闷而更加感到倦怠。}。莫离散而发曙兮\endnote{莫离散：指晚潮退去。莫，通“暮”。发曙：指早潮发生而来。}，内存心而自持\endnote{内：指内心。存心而自持：保存着观涛的印象自可持久不忘。}。于是澡概胸中\endnote{澡概：同“澡溉”，洗涤。此句言观涛之后心中如同经过洗涤似的。}，洒练五藏\endnote{洒：洗。练：汰。五藏：即五脏。}，澹澉手足\endnote{澹澉（hàn旱）：洗涤。}，颒濯发齿\endnote{颒（huì汇）濯：洗涤。颒，指洗脸。}。揄弃恬怠\endnote{揄（shū书）：通“输”，脱，抛弃。《广雅·释诂四》：“揄，脱也。”恬怠：懒散，倦怠。}，输写淟浊\endnote{输写：同“输泻”，排除。淟浊：污浊。指身体中的污垢杂质。}，分决狐疑\endnote{分决：判明决定。狐疑：犹豫不定，怀疑。}，发皇耳目\endnote{发皇耳目：使耳目受到启发而耳聪目明。皇，明的意思。}。当是之时，虽有淹病滞疾\endnote{淹病滞疾：沉滞日久的疾病。}，犹将伸伛起躄\endnote{伸伛（yǔ羽）：使伛偻的人伸直了腰。起躄（bì毕）：使足不能走路的人可以站起来行走。躄，足不能行，即中医所说的“痿躄”。}，发瞽披聋而观望之也\endnote{发瞽：使盲人睁开眼睛。发，开。瞽，盲人。披聋：使聋子听得见声音。披，劈开，打开。之：指代江涛。}，况直眇小烦懑、酲醲病酒之徒哉\endnote{况直：何况只是。眇小烦懑：指小病与不舒服。酲醲病酒：指饮酒过量引起的疾病。}？故曰发蒙解惑，不足以言也\endnote{“发蒙”二句：语本《黄帝内经·素问》：“发蒙解惑，未足以论也。”言使昏惑的头脑清楚，就不值得一提了。}。”太子曰：“善，然则涛何气哉\endnote{何气：何种气象。}？”

客曰：“不记也\endnote{不记：指不见记载。}。然闻于师曰，似神而非者三\endnote{“似神”句：言江涛若有神助而又并非神力所致的特点有三个。}：疾雷闻百里\endnote{“疾雷”句：言江涛如迅疾的雷声，声闻百里之遥。这是特点之一。}；江水逆流，海水上潮\endnote{“江水”二句：言海水涨潮时，长江之水逆流而上。这是特点之二。}；山出内云，日夜不止\endnote{“山出内云”二句：言江涛雾气笼罩如同山峰吞吐云气一样，而且日夜不停。这是第三个特点。内，同“纳”。}。衍溢漂疾\endnote{衍溢：指大水满溢。漂疾：指水流迅速。}，波涌而涛起。其始起也，洪淋淋焉\endnote{洪：指浩荡的潮水洪峰。淋淋：水自上而下洒落的样子。《说文》：“淋淋，山下水貌。”}，若白鹭之下翔。其少进也\endnote{少进：对“始起”而言，指稍有进展。}，浩浩溰溰\endnote{浩浩溰（qì气）溰：浩大洁白的样子。}，如素车白马帷盖之张\endnote{“如素车”句：言潮水如高张帷盖的白色车马。}。其波涌而云乱\endnote{云乱：形容潮水如乱云翻滚。}，扰扰焉如三军之腾装\endnote{扰扰：纷乱的样子。腾装：整理装束。六臣注《文选》刘良注此二句说：“言如云气之乱，又如三军之装束也。”}。其旁作而奔起也\endnote{旁作：指潮头横出。奔起：指潮头上扬。}，飘飘然如轻车之勒兵\endnote{如轻车之勒兵：如同乘坐轻便的战车指挥军队。轻车，古代兵车名，为战车中最轻便者。勒：部署；指挥。}。六驾蛟龙，附从太白\endnote{“六驾”二句：言六匹驾车的蛟龙跟随河神之后而受其指挥。这是借以形容江潮之势。太白，即河伯（用李善说）。一说太白指帅旗（张凤翼说），亦可通。}。纯驰浩霓\endnote{纯驰：纯通“屯”，集的意思。指浪涛积聚成峰，若屯驻不行。驰，指涛之奔驰。浩：通“皓”，白。皓霓，即白虹。用以喻涛。此句写江涛或屯聚成峰，或奔驰若白虹行天。}，前后骆驿\endnote{骆驿：同“络绎”。前后连续不断。}。颙颙卬卬\endnote{颙（yóng喁）颙卬卬：浪涛高大的样子。}，椐椐强强\endnote{椐（jú局）椐强强：形容浪涛横恣之态（用吕锦文《文选古字通训补》说）。李善注谓“相随之貌”。}，莘莘将将\endnote{莘（shēn伸）莘将将：浪涛互相激荡的样子（用李周翰说）。}。壁垒重坚，沓杂似军行\endnote{“壁垒”二句：言浪涛如军营的壁垒，重叠而坚固，其盛大众多，如同军队的行列。沓杂，同“杂沓”，五臣本《文选》作“杂沓”，纷杂繁多的样子。}。訇隐匈磕\endnote{訇（hōng轰）隐匈磕（kē科）：象声词，皆形容江涛所发出的巨大声音。}，轧盘涌裔\endnote{轧盘：广大的样子。轧，即坱轧无垠。盘，即盘礴。二字并言广大。涌裔：波浪沸腾翻滚的样子。此句言涛势广大，沸腾滚动。}，原不可当\endnote{原不可当：实不可抵挡。原，本。}。观其两傍，则滂渤怫郁\endnote{滂渤怫（fú拂）郁：水势怒激的样子。}，暗漠感突\endnote{暗漠：本义为昏暗不明。此指江涛烟波浩渺，看不真切。感突：指波浪左右冲突。感，通“撼”。}。上击下律\endnote{上击：指潮头上升如物击向空中。下律：律，当作“硉”（lǜ律），推石而下叫硉，下律言潮头从半空而下如推石坠落。}，有似勇壮之卒，突怒而无畏\endnote{突怒：左右奔突而发怒。}。蹈壁冲津\endnote{蹈壁：指江涛拍打岸壁。蹈，本义为踏，此用引申义。冲津：冲击渡口。}，穷曲随隈\endnote{穷曲随隈：江湾深曲隈隩之处，江涛无所不至。穷，穷尽。随，追随。隈，江湾深曲之处。}，逾岸出追\endnote{逾岸：越过江岸。出追：超出沙堆。追，古“堆”字。}，遇者死，当者坏。初发乎或围之津涯\endnote{初发：指潮水最初的发源地。或围：李善注谓地名，不可考，或为虚拟之地名。津涯：水涯渡口。}，荄轸谷分\endnote{荄轸：荄，通“陔”，作陇解，指山陇。轸，转动。谷分：山谷为之分裂。此句言广陵潮所到之处，山陇为之转动，山谷为之分裂。}，回翔青篾\endnote{回翔青篾：言涛在青篾这个地方回旋。青篾，地名（用李善注）。一说青篾为车名（胡绍煐说）。}，衔枚檀桓\endnote{衔枚：横衔枚于口中，以防止喧哗或叫喊。枚，形如筷子，两端有带，可系于颈上。此处用衔枚喻潮水突然无声，如衔枚似的。檀桓：地名（用李善注）。}，弭节伍子之山\endnote{弭节：注见〔187〕。伍子之山：当指胥山，在太湖边，离长江不过百里。按：《史记·伍子胥列传》说：伍子胥临死之时，曾告诉其舍人说：“抉吾眼县（悬）吴东门之上，以观越寇之入灭吴也。”乃自刎而死。吴王闻之大怒，乃取子胥尸盛以鸱夷革，浮之江中。吴人怜之，为立祠于江上，因命曰胥山。又《论衡·书虚》篇说：“吴王杀子胥，投之江，子胥恚恨，驱水为涛，以溺杀人。”在汉代已有这种传说，一说子胥死后，化为怒涛，故凡有江涛之处，就可能有纪念伍子胥的祠庙，江都县有江水祠，俗谓之伍相庙，所以伍子之山，亦可能不是确指。}，通厉骨母之场\endnote{通厉：远行。骨母：当为“胥母”之误。李善注引《史记》说：“吴王杀子胥，投之于江，吴人立祠于江上，因名胥母山。”前注引《史记·伍子胥列传》原文，称“胥山”，无“母”字，或指同一山。一说胥山即吴山，在浙江杭州。胥母山在江苏吴县西南，主广陵潮即钱塘潮说者往往以此为据。汪中《述学·广陵曲江证》已驳此说，认为伍子之山、胥母之场皆与浙江无关。案：场乃平坦空地，疑胥母之场，乃言江涛离开胥母山后进入平坦空地又向远方流去。即便伍子之山与胥母之山为同一山，亦不相矛盾。}。凌赤岸\endnote{凌：凌驾，凌跨。指涛言。赤岸：地名，在广陵（用李善注）。}，彗扶桑\endnote{彗：作动词用，扫、拂之意。扶桑：指东方日出之处。}，横奔似雷行\endnote{“横奔”句：言江涛横流如迅雷奔行。}。诚奋厥武\endnote{诚：确实。奋厥武：奋发了江涛的威武。厥，指代江涛。}，如振如怒\endnote{如振如怒：像响雷像愤怒。振，同“震”，响雷。}。沌沌浑浑\endnote{沌沌浑浑：波涛相随的样子（用李善注）。}，状如奔马。混混庉庉\endnote{混混庉（tún屯）庉：指波浪的声音（用李善注）。}，声如雷鼓\endnote{声如雷鼓：声音像雷鸣，又像鼓响。}。发怒庢沓\endnote{发怒庢（zhì至）沓：言江涛怒不可遏，遇到阻碍，则从旁涌溢而出。庢，阻碍。沓，李善注引《埤苍》：“釜沸出也。”即沸水从锅中涌出。}，清升逾跇\endnote{清升：指涛平缓时清波上扬。逾跇（yì意）：指水浪跳跃着互相超越。}，侯波奋振\endnote{侯波：指大波。侯，指阳侯，传说为波涛之神，亦借指波涛。奋振：奋力振起。}，合战于藉藉之口\endnote{合战：会战，相斗。藉藉之口：名叫藉藉的港口。藉藉，虚拟的地名。}。鸟不及飞，鱼不及回，兽不及走。纷纷翼翼\endnote{纷纷翼翼：众多波涛交错的样子。}，波涌云乱。荡取南山\endnote{荡取：荡，冲激。取，同“趣”，即“趋”，奔向。}，背击北岸\endnote{背击北岸：言从南山反回之波又撞击江之北岸。背，反。}。覆亏丘陵\endnote{“覆亏”句：言江涛可使丘陵倾覆、损坏。}，平夷西畔\endnote{“平夷”句：言潮水与两岸相平。平夷，平坦。畔，岸。}。险险戏戏\endnote{险险戏戏：崎岖危险的样子。戏，同“巇”（xī西）。此句亦指涛而言。}，崩坏陂池\endnote{“崩坏”句：言江涛把坡形的堤岸都冲坏了。陂池，通“陂陁”，本指斜坡，这里指堤岸。}，决胜乃罢\endnote{“决胜”句：言江涛会战取得胜利之后才罢休。呼应前文“合战于藉藉之口”句。}。濞汩潺湲\endnote{濞汩（zhì yù制玉）：水流激荡的样子。潺湲：水流之声。此句言水波相激发出潺潺之声。}，披扬流洒\endnote{披扬：指波涛翻腾。流洒：指浪花飞溅。}，横暴之极。鱼鳖失势\endnote{失势：指失去常态。}，颠倒偃侧\endnote{颠倒偃侧：指鱼鳖被江涛掀翻了个儿。偃，仰，此指颠倒为肚皮朝上。侧，侧身，指翻了个半身。}。沋沋湲湲\endnote{沋（yóu油）沋湲湲：李善注谓“鱼鳖颠倒之貌”。吕延济注：“沋沋湲湲，皆声也。”按：二家注均缺乏训诂根据，乃以意测之，李说近是。}，蒲伏连延\endnote{蒲伏：通“匍匐”，指鱼鳖伏地爬行。连延：连绵不断。}。神物怪疑，不可胜言。直使人踣焉\endnote{直使人踣（bó勃）：简直使人受惊倒下。踣，向前仆倒。}，洄暗凄怆焉\endnote{洄暗：昏瞆不明的样子。五臣注《文选》刘良注：“洄暗，深不明也。”一说：“惊骇失智貌。”}。此天下怪异诡观也\endnote{怪异诡观：指怪异之物与诡奇的景观。}。太子能强起观之乎？”太子曰：“仆病，未能也。”

客曰：“将为太子奏方术之士\endnote{奏：进，献。方术之士：此指精通学术的人。实指战国诸子。}，有资略者\endnote{资略：指才资谋略。}，若庄周、魏牟、杨朱、墨翟、便蜎、詹何之伦\endnote{庄周：庄子，名周，战国时道家学派的代表人物。魏牟：战国时人，又称“中山公子牟，魏公子牟”。《吕氏春秋·审为篇》载：“中山公子牟谓詹子曰：‘身在江海之上，心居魏阙之下，奈何？’詹子曰：‘重生，重生则轻利。’中山公子牟曰：‘虽知之，犹不能自胜也。’詹子曰：‘不能自胜则纵之，神无恶乎。不能自胜而强不纵者，此之谓重伤，重伤之人无寿类矣。’”高诱注：“子牟，魏公子也，作书四篇。魏伐得中山，公以邑子牟，因曰中山公子牟也。”杨朱：战国时人，诸子之一，称杨子，他是个利己主义者，主张拔一毛利于天下而不为。（见《孟子·尽心下》）墨翟：即墨子，战国时的思想家，是墨家学派的代表人物，著有《墨子》一书。便蜎（juān娟）：即蜎蠉，白公（即楚国之白公胜）时人。也即蜎子。《七略》说：“蜎子，名渊，楚人也。”《两汉文学史参考资料》引林茂春说，以为便蜎即环渊。按：环渊之名，一见于《史记·田完敬仲世家》：“宣王喜文学游说之士，自如驺衍、淳于髡、田骈、接予、慎到、环渊之徒七十六人，皆赐列第，为上大夫，不治而议论。”据此知环渊为稷下学者之一。《史记正义》注说：“楚人。《孟子传》云环渊注书上下篇也。”又见于《史记·孟子荀卿列传》，称环渊为“稷下先生”。《史记索隐》按语说：“刘向《别录》‘环’作姓也。”又《汉书·艺文志》著录《蜎子》十三篇，今已不存。詹何：即詹子，战国时代人，《庄子》、《吕氏春秋》、《列子》等书都提到此人。《吕氏春秋·审为篇》载中山公子牟与詹子的对话，前已引。高诱注说：“詹子，古得道者也。”又《吕氏春秋·执一篇》载“楚王问为国于詹子”，高诱注：“詹何，隐者。”陈奇猷按语说：“‘詹’，《庄子·让王》作‘瞻’，字通。《列子·汤问》张湛注：‘詹何，楚人。……考《庄子·让王》篇及本书《审为篇》载詹何与中山公子牟答问，则詹何当是楚顷襄王时人（《汉书·古今人表》列公子牟与楚顷襄王同时）。……《韩非子·解老》云：“詹何坐，有牛鸣于门外，詹何曰：‘是黑牛也，而白在其角。’使人视之，是黑牛而以布裹其角。”《淮南·原道训》云：“詹何娟嬛之数”，又《览冥训》云：“詹何骛鱼大渊之中”（事详《列子·汤问》），则詹何盖一术士。《淮南·诠言训》引有詹何语。本书《先己》、《孝行》、《求人》、《处方》四篇亦皆引詹何曰，似詹何本有著述，而《汉书·艺文志》失载。’”}，使之论天下之精微\endnote{精微：精辟微妙的道理。按：精微，底本作“释微”。《文选·胡氏考异》云：“五臣本作‘精’。……（李）善引‘好论精微’为注，似亦作‘精’。各本所见，皆传写误。”据此改“释”为“精”。}，理万物之是非\endnote{理：治理，申辨。}，孔、老观览\endnote{孔、老：指孔子和老子。观览：这里有裁判、评断之义。}，孟子持筹而算之\endnote{持筹而算之：手拿计算工具以记数。筹，指古代刻有数字的竹筹计算工具。算，记数，此指记下各人胜败之数。}，万不失一。此亦天下要言妙道也。太子岂欲闻之乎？”于是太子据几而起\endnote{据：凭依。几：古人坐时凭依或搁置物件的小桌。}，曰\endnote{曰：当为衍文。吴闿生说：“此‘曰’字当是衍文，删去之，文乃可读。”（见《古文辞类纂音注》引）按：后一句，不像太子的语言，吴说是。}：“涣乎若一听圣人辩士之言\endnote{“涣乎”句：言太子精神顿时清醒，如同听到圣人辩士的话。涣，清醒，明白。}。”涊然汗出\endnote{涊（niǎn捻）然：汗出的样子。}，霍然病已\endnote{霍然：犹“忽然”。按：《玉篇·雨部》：“霍，鸟飞急疾貌也。”已：止。}。

\printendnotes